\documentclass[11pt,a4paper]{article}
\usepackage[T1]{fontenc}
\usepackage[utf8]{inputenc}
\usepackage[spanish,es-nodecimaldot,es-noshorthands]{babel}
\usepackage{lmodern,microtype,amsmath,amssymb,siunitx}
\usepackage[margin=23mm]{geometry}
\setlength{\parindent}{0pt}\setlength{\parskip}{5pt}
\setlength{\emergencystretch}{2em}
\newcommand{\dd}{\mathrm{d}}
\newcommand{\PadSide}{0.5}
\newcommand{\PadHalf}{0.25}
\newcommand{\GapSize}{5}
\newcommand{\GapSpeed}{3}
\newcommand{\GapField}{6}
\newcommand{\Mobility}{5}
\newcommand{\PeakCurrent}{1.731}
\newcommand{\GapTime}{1.667}
\newcommand{\GapPhi}{0.0015876}
\newcommand{\GapFraction}{0.9984124}
\newcommand{\FinalTime}{193.5}
\newcommand{\FrameTime}{10}
\newcommand{\Plateau}{0.016}
\newcommand{\ExampleCharge}{1.6}
\newcommand{\ExampleCurrent}{12.82}
\newcommand{\DerivativeError}{9.75e-11}
\newcommand{\ChargeError}{1.51e-14}

\title{Corriente inducida en un pad de NAUSICA}
\author{J. J. G\'omez-Cadenas\\{\small DIPC}}
\date{\today}
\begin{document}
\maketitle
Rogelio, la derivación de la corriente está en la sección
\emph{Signal formation and response of one pad} del artículo
\emph{NAUSICAA for ITACA}. El apéndice \emph{Analytical single-ion response}
recoge las condiciones de contorno, el dibujo del potencial de ponderación
(\emph{weighting potential}) y las comprobaciones del cálculo.

En la expresión del pico de corriente, en \emph{Signal size and duration},
$a$ es la \textbf{mitad del lado del pad}. El pad es cuadrado y su lado mide
$2a=\SI{\PadSide}{mm}$, por tanto
\[
 \boxed{a=\SI{\PadHalf}{mm}=\SI{2.5e-4}{m}.}
\]
El pico de \SI{\PeakCurrent}{fA} corresponde a \textbf{un ion de carga $+e$}
justo antes de llegar al centro del pad, con la velocidad de
\SI{\GapSpeed}{m/s} adoptada en el artículo.

\section*{Comparación con la cuenta de tu esquema}
La expresión $I=Qv/d$ corresponde a un campo de ponderación uniforme,
$W=1/d$, como el de dos placas planas paralelas separadas por $d$.
En tu cálculo asignas al pad la carga total de unos $10^5$ iones.
La geometría de condensador utilizada y la asignación de esa carga total
al pad no representan la señal de entrada de un canal de NAUSICA.
Por tanto, la respuesta del CSA a esa excitación no permite evaluar su
sensibilidad a los iones que llegan a un pad del detector.

Para calcular esa señal necesitamos el campo de ponderación del pad y
los tiempos de llegada de los iones que recoge ese canal. El campo de
ponderación aumenta cerca del pad, y cada ion produce una corriente que
varía durante su aproximación. La sección \emph{Response to uniform arrivals}
del artículo explica cómo construimos la entrada de SPICE sumando estas
corrientes para el caso de llegadas uniformes durante un frame de
\SI{\FrameTime}{ms}. El pico de \SI{\PeakCurrent}{fA} de un ion no se
mantiene durante todo el frame.

Por separado, hay una inconsistencia en la anotación de carga: los
\SI{16}{fC} utilizados corresponden aproximadamente a $10^5$ iones,
es decir, $100\,\mathrm{ke}$. La etiqueta «$10\,\mathrm{ke}$» corresponde
a $10^4$ cargas elementales (\SI{\ExampleCharge}{fC}) y debería corregirse
a «$100\,\mathrm{ke}$».
\end{document}
